\writeSectionFrame{\presentationTitle}{Das Strategie-Muster}
\begin{frame}{Steckbrief}
	\begin{itemize}
 		\item Deutscher Name: Strategie
 		\item Alternativer deutscher Name: Policy
 		\item Englischer Name: Strategy
 		\item Gruppe: Verhaltensmuster
 	\end{itemize}
\end{frame}

\begin{frame}{Strategie}
	\begin{block}{Zweck}
 		Definiere eine Familie von Algorithmen, kapsele jeden einzelnen und mache sie austauschbar. Der Algorithmus kann unabhängig von den nutzenden Klienten variiert werden.
 	\end{block}
\end{frame}

\begin{frame}{Klassendiagramm}
\begin{tikzpicture}[scale=0.6, transform shape]

    \umlclass{Context}{
        -strategy: Strategy
    }{
        +Context(Strategy)\\
        +setStrategy(Strategy)\\
        +executeStrategy(): void\\
    }

    \umlclass[y=-3.5]{Strategy}{}{
        +execute(): void
    }

    \umlclass[x=-4, y=-6]{ConcreteStrategyA}{}{
        +execute(): void
    }

    \umlclass[x=4, y=-6]{ConcreteStrategyB}{}{
        +execute(): void
    }

    \umlinherit{ConcreteStrategyA}{Strategy}
    \umlinherit{ConcreteStrategyB}{Strategy}
    \umlaggreg{Context}{Strategy}

\end{tikzpicture}
\end{frame}

\note{
	Strategie definiert eine einheitliche Schnittstelle zu allen unterstützten Algorithmen.
	Kontext benutzt diese Schnittstelle um einen konkreten Algorithmus aufzurufen. KonkreteStrategie
	implementiert den Algorithmus mit Hilfe der Strategie-Schnittstelle. Kontext wird mit einem 
	KonkreteStrategie-Objekt konfiguriert. Es behält eine Referenz zu einem Strategie-Objekt und 
	kann eine Schnittstelle zur Verfügung stellen, um der Strategie Daten zu übergeben. Kontext 
	kann natürlich auch Referenzen auf mehrere Strategieobjekte haben. Nach dem Erzeugen des 
	Kontextes und eventuell einer Strategie (der Kontext selbst kann die Strategie aber auch erzeugen)
	interagiert der Klient nur noch mit dem Kontext. 
}

\begin{frame}{Klassendiagramm - Beispiel}
\begin{tikzpicture}[scale=0.6, transform shape]
    \umlclass{TransportService}{
        -strategy: TransportStrategy
    }{
        +TransportService(TransportStrategy) \\
        +setStrategy(TransportStrategy)\\
        +calculateTransportCost(distance: double): double\\
    }

    % Strategie-Interface
    \umlclass[y=-3]{TransportStrategy}{}{
        +calculateCost(distance: double): double
    }

    % Konkrete Strategien
    \umlclass[x=-4, y=-6]{CarTransportStrategy}{}{
        +calculateCost(distance: double): double
    }

    \umlclass[x=4, y=-6]{TrainTransportStrategy}{}{
        +calculateCost(distance: double): double
    }

    \umlclass[x=0, y=-9]{AirplaneTransportStrategy}{}{
        +calculateCost(distance: double): double
    }

    % Beziehungen
    \umlinherit{CarTransportStrategy}{TransportStrategy}
    \umlinherit{TrainTransportStrategy}{TransportStrategy}
    \umlinherit{AirplaneTransportStrategy}{TransportStrategy}
    \umlaggreg{TransportService}{TransportStrategy}

\end{tikzpicture}
\end{frame}

\begin{frame}[fragile]{Strategie}
	\begin{exampleblock}{\javapackage{sort.SortingStrategy}}
\begin{lstlisting}
public interface SortingStrategy {
	public void sort(int[] toSort);
}
\end{lstlisting}
 	\end{exampleblock}
\end{frame}

\begin{frame}[fragile]{Konkrete Strategie}
	\begin{exampleblock}{\javapackage{sort.BubbleSort}}
 		\begin{lstlisting}
public class BubbleSort implements SortingStrategy {
	@Override
	public void sort(int[] toSort) {
		for (int i = 1; i < toSort.length; i++) {
			for (int j = 0; j < toSort.length - i; j++) {
				if (toSort[j] > toSort[j + 1]) {
					int temp = toSort[j];
					toSort[j] = toSort[j + 1];
					toSort[j + 1] = temp;
				}
			}
		}
	}
}
 		\end{lstlisting}
 	\end{exampleblock}
\end{frame}

\begin{frame}[fragile]{Kontext}
	\begin{exampleblock}{\javapackage{sort.SortingContext}}
 		\begin{lstlisting}
public class SortingContext {
	private static final int STRATEGY_COUNT = 3;
	private SortingStrategy[] strategies;
	private int[] toSort;
	
	public SortingContext(int[] zusortieren) {
		strategies = new SortingStrategy[STRATEGY_COUNT];
		strategies[0] = new QuickSort();
		strategies[1] = new InsertionSort();
		strategies[2] = new BubbleSort();
		
		this.toSort = zusortieren;
	}
 		\end{lstlisting}
 	\end{exampleblock}
\end{frame}

\begin{frame}[fragile]{Kontext}
	\begin{exampleblock}{\javapackage{sort.SortingContext}}
 		\begin{lstlisting}
public TimeMeasurement[] sort() {
    TimeMeasurement[] measurements = new 
        TimeMeasurement[STRATEGY_COUNT];
    
    for (int i = 0; i < STRATEGY_COUNT; i++) {
        SortingStrategy sortingStrategy = strategies[i];
        
        int[] copyOfArray = Arrays.copyOf(toSort, 
            toSort.length);
        final long start = System.currentTimeMillis();
        sortingStrategy.sort(copyOfArray);
        final long end = System.currentTimeMillis();
        measurements[i] = new TimeMeasurement(start, end, 
            sortingStrategy);
    }
    
    return measurements;
}
 		\end{lstlisting}
 	\end{exampleblock}
\end{frame}

\begin{frame}[fragile]{Klient}
	\begin{exampleblock}{\javapackage{dialog.SortingDialog}}
 		\begin{lstlisting}
private int[] toSort;
private SortingContext sortingContext;

public SortingDialog() {
    toSort = new int[VALUE_COUNT];
    fillArray();
}

public void start() {
    sortingContext = new SortingContext(toSort);
    TimeMeasurement[] timeMeasurements = sortingContext.sort();
    
    for (TimeMeasurement timeMeasurement: timeMeasurements) {
        System.out.println(timeMeasurement.toString());
    }
}	    
 		\end{lstlisting}
 	\end{exampleblock}
\end{frame}

\begin{frame}{Klassendiagramm - Beispiel aus dem JDK}
	\begin{center}
 		\includegraphics[scale=0.5]{\BILDERPFAD/strategie3.png}
 	\end{center}
\end{frame}

\begin{frame}{Klassendiagramm - Zahlungssysteme}
	\begin{center}
 		\includegraphics[scale=0.4]{\BILDERPFAD/strategie4.png}
 	\end{center}
\end{frame}

\begin{frame}[fragile]{Strategie}
	\begin{exampleblock}{\javapackage{payment.PaymentStrategy}}
 		\begin{lstlisting}
public interface PaymentStrategy {
	void pay(Person person, int amount);
}
 		\end{lstlisting}
 	\end{exampleblock}
\end{frame}

\begin{frame}[fragile]{Konkrete Strategie}
	\begin{exampleblock}{\javapackage{PayPalStrategy}}
 		\begin{lstlisting}
public class PayPalStrategy implements PaymentStrategy {
	@Override
	public void pay(Person person, int amount) {
		PayPalData payPalData = person.getPayPalData();
		
		if (payPalData == null) {
			throw new IllegalArgumentException("...");
		}
		
		if (!payPalData.getPassword().equals("geheim")) {
			throw new IllegalArgumentException("...");
		}
		
		person.payWithPayPal(amount);
	}
} 		    
 		\end{lstlisting}
 	\end{exampleblock}
\end{frame}

\begin{frame}[fragile]{Kontext}
	\begin{exampleblock}{\javapackage{payments.OnlineShop}}
        \begin{lstlisting}
public class OnlineShop {
	public void pay(Payments payment, Person person, int amount) {
		PaymentStrategy paymentStrategy = null;
		
		switch (payment) {
		case CREDITCARD: paymentStrategy = 
            new CreditCardStrategy();
		break;
		case DEBITENTRY: paymentStrategy = 
            new DebitEntryStrategy();
		break;
		case PAYPAL: paymentStrategy = 
            new PayPalStrategy();
		break;
		}
		paymentStrategy.pay(person, amount);
	}
}    
        \end{lstlisting}
 	\end{exampleblock}
\end{frame}

\begin{frame}[fragile]{Klient}
	\begin{exampleblock}{\javapackage{dialog.PaymentDialog}}
        \begin{lstlisting}
private int chooseFunction(int amount) {
    OnlineShop onlineShop = new OnlineShop();
    
    switch (menuItem) {
    case CREDITCARD: onlineShop.pay(
        Payments.CREDITCARD, person, amount);
    break;
    case PAYPAL: onlineShop.pay(
        Payments.PAYPAL, person, amount);
    break;
    case DEBITENTRY: onlineShop.pay(
        Payments.DEBITENTRY, person, amount);
    break;
    }
    
    return function;
} 
        \end{lstlisting}
 	\end{exampleblock}
\end{frame}

\begin{frame}{Klassendiagramm - Speichersysteme}
	\begin{center}
 		\includegraphics[scale=0.4]{\BILDERPFAD/strategie5.png}
 	\end{center}
\end{frame}

\begin{frame}[fragile]{Strategie}
	\begin{exampleblock}{\javapackage{queue.Queue}}
 		\begin{lstlisting}
public interface Queue {
	void addNumber(int number);
	int getNumber(int index);
	int popNumber();
	int peekNumber();
	int size();
}
 		\end{lstlisting}
 	\end{exampleblock}
\end{frame}

\begin{frame}[fragile]{Konkrete Strategie}
	\begin{exampleblock}{\javapackage{queue.ArrayQueue}}
 		\begin{lstlisting}
public class ArrayQueue implements Queue {
	private Integer[] numbers;
	private int count;
	
	public ArrayQueue() {
		numbers = new Integer[SIZE];
		count = 0;
	}
	
	@Override
	public void addNumber(int number) {
		if (count < numbers.length) {
			numbers[count] = number;
			count++;
		} else {
			throw new ArrayFullException();
		}
	}   
 		\end{lstlisting}
 	\end{exampleblock}
\end{frame}

\begin{frame}[fragile]{Kontext}
	\begin{exampleblock}{\javapackage{queue.QueueContext}}
        \begin{lstlisting}
public class QueueContext {
	private Queue queue;
	
	public QueueContext(boolean duplicatesAllowed, int count) {
		if (!duplicatesAllowed) {
			queue = new SetQueue();
			System.out.println("Eine SetQueue wurde erzeugt");
		} else if (count <= 5) {
			queue = new ArrayQueue();
			System.out.println("Eine ArrayQueue wurde erzeugt");
		} else {
			queue = new ListQueue();
			System.out.println("Eine ListQueue wurde erzeugt");
		}
	}
        \end{lstlisting}
 	\end{exampleblock}
\end{frame}

\begin{frame}[fragile]{Kontext}
	\begin{exampleblock}{\javapackage{queue.QueueContext}}
        \begin{lstlisting}
public int peekNumber() {
    return queue.peekNumber();
}

public int popNumber() {
    return queue.popNumber();
}

@Override
public String toString() {
    return queue.toString();
}
        \end{lstlisting}
 	\end{exampleblock}
\end{frame}

\begin{frame}[fragile]{Kontext}
	\begin{exampleblock}{\javapackage{queue.QueueContext}}
        \begin{lstlisting}
public void addNumber(int number) {
    try {
        queue.addNumber(number);
    } catch (ArrayFullException ex) {
        Queue arrayQueue = queue;
        System.out.println("Wir schalten um auf ...");
        queue = new ListQueue();
        
        for (int i = 0; i < arrayQueue.size(); i++) {
            int oldNumber = arrayQueue.getNumber(i);
            queue.addNumber(oldNumber);
        }
        
        queue.addNumber(number);
    }
}
        \end{lstlisting}
 	\end{exampleblock}
\end{frame}

\begin{frame}[fragile]{Klient}
	\begin{exampleblock}{\javapackage{dialog.QueueDialog}}
        \begin{lstlisting}
private void addValue() {
    System.out.println("Bitte eine Zahl eingeben: ");
    int number = scanner.nextInt();
    queueContext.addNumber(number);
}

private void createNewQueue() {
    System.out.println("Wie viele Zahlen sollen eingegeben " +
        " werden?");
    int count = scanner.nextInt();
    System.out.println("Sollen doppelte Werte erlaubt sein? " + 
        true/false");
    boolean duplicatesAllowed = scanner.nextBoolean();
    queueContext = new QueueContext(duplicatesAllowed, count);
}
        \end{lstlisting}
 	\end{exampleblock}
\end{frame}

\begin{frame}[fragile]{Klient}
	\begin{exampleblock}{\javapackage{dialog.QueueDialog}}
        \begin{lstlisting}
private int chooseFunction() {
    switch (menuItem) {
    case NEW_QUEUE: createNewQueue();
    break;
    case ADD: addValue();
    break;
    case PEEK: System.out.println("Erster Wert: " + 
        queueContext.peekNumber());
    break;
    case POP: System.out.println("Entfernt: " +     
        queueContext.popNumber());
    break;
    case SHOW: System.out.println(queueContext.toString());
    break;
    }
    
    return function;
}
        \end{lstlisting}
 	\end{exampleblock}
\end{frame}